%!TEX root = TheoretischePhysikC.tex
\section{Vorlesung: Unschärfe}
Wir betrachten nun freie Teilchen, die im Hilbertraum $\mathcal{H}\equiv L^2(\mathbb{R}^3)$ beschrieben werden. Der Raum $\mathbb{R}^3$ ist hier als Ortsraum interpretiert, das heißt, in dieser Darstellung des Hilbertraums sind die Funktionen der Ortsobservable als Multiplikationsoperatoren dargestellt. In einem reinen Zustand $\rho = |\Psi\rangle\langle \Psi|$ ist daher die Wahrscheinlichkeitsdichte für die Ergebnisse von Ortsmessungen $p(\mathbf{x})\equiv |\Psi(\mathbf{x})|^2$.

Eine Ortsmessung ist dabei ein Gitter von Ja/Nein-Detektoren, im Grunde eine Kamera; der Operator $\widehat{x}_j$ selbst wird nie direkt gemessen, er steht für dieses Verfahren.

Der Erwartungswert der Ortskomponente $\widehat{x}_j$ ergibt sich durch Integration:
\begin{equation}
	\mathbb{E}[x_j] = \iiint x_j |\Psi(\mathbf{x})|^2\, dxdydz = \tr(\rho \widehat{x}_j),
\end{equation}
und entsprechend für höhere Momente. Speziell ist also die Varianz der Verteilung von $x_j$:
\begin{equation}
	\text{var}(x_j) = \tr(\rho \widehat{x}_j^2) - \tr(\rho \widehat{x}_j)^2.
\end{equation}

Für den Impulsoperator erhalten wir analog
\begin{align}
	\mathbb{E}[p_j] &= \tr(\rho \widehat{p}_j) \\
	\text{var}(p_j) &= \tr(\rho \widehat{p}_j^2) - \tr(\rho \widehat{p}_j)^2
\end{align}

Das sind zwei getrennte Experimente: Entweder misst man den Ort, oder man misst den Impuls. Bis hierher unterscheiden sich die Vorhersagen von klassischer Physik und Quantenmechanik nicht; der Bruch kommt jetzt.

Es stellte sich heraus, dass nicht beide Verteilungen zugleich scharf konzentriert sein können. Das heißt: beide Varianzen können nicht beliebig klein werden. Seien $A$ und $B$ beliebige hermitesche Operatoren. Dann ist für alle 
$\lambda \in \mathbb{R}$ der Operator $M = (A+i\lambda B)^\dag (A+i\lambda B)$ (per definition) positiv,
hat also positive Erwartungswerte. Dass $\lambda$ reell sein muss, ist wesentlich: Die Adjungierung
konjugiert komplexe Zahlen, und für komplexes $\lambda$ ginge die Rechnung schief. Der Operator $M$
ist ein Kunstgriff; entscheidend ist allein, dass er die Form $C^\dag C$ mit $C=A+i\lambda B$ hat, denn diese Form garantiert Positivität. Damit ist für $\lambda\in \mathbb{R}$ und eine beliebige
Dichtematrix $\rho$:
\begin{equation}
	\tr(\rho M) = \tr(\rho A^2)+\lambda^2\tr(\rho B^2) + \lambda\tr(\rho i[A,B])\ge 0.
\end{equation}
Setzt man hier für $\lambda$ den Wert ein, bei dem die linke Seite am kleinsten wird (Übung),
bekommt man:
\begin{equation}
	\tr(\rho A^2)\tr(\rho B^2) \ge \frac14 \tr(\rho i[A,B])^2.
\end{equation}
Das Rezept dafür: nach $\lambda$ ableiten, die Ableitung null setzen. Die Ungleichung gilt für beliebige hermitesche Operatoren.
Setzt man hier ein: $A=\widehat{p}_j - \tr(\rho \widehat{p}_j )\mathbb{I}$ und $B=\widehat{x}_j - \tr(\rho \widehat{x}_j )\mathbb{I}$, und benutzt die kanonischen Vertauschungsrelationen für die rechte Seite, so bekommt man
\begin{equation}
	\eqbox{\text{var}(p_j)\text{var}(x_j) \ge \frac{\hbar^2}{4}, \quad\text{oder}\quad (\Delta p_j)(\Delta x_j) \ge \frac{\hbar}{2},}
\end{equation}
wenn man die $(\Delta A)$ als Wurzel der Varianz (Standardabweichung) einführt. Dies ist die \emph{Heisenberg'sche Unschärferelation}, die in dieser Form von Kennard stammt.

Die Interpretation der Heisenberg'schen Unschärferelation lautet: Es ist unmöglich, quantenmechanische Teilchen so zu präparieren, dass die Varianzen bei Impuls-und Ortsmessungen beide beliebig klein ausfallen. Die Größen $(\Delta p_j)$ und $(\Delta x_j)$ beziehen sich also auf völlig verschiedene statistische Experimente. Eine Standardabweichung entsteht erst durch viele Wiederholungen: präparieren, messen, wieder genauso präparieren, wieder messen; $(\Delta p_j)$ bestimmt man in einem davon getrennten Durchgang.

\begin{merksatz}[Was die Unschärferelation nicht sagt]
Sie sagt nichts über sukzessive Messungen. Die verbreitete Darstellung, man messe erst den Ort und dann am selben Teilchen den Impuls, ist falsch, auch wenn fast alle Autoren sie so bringen. Für Messfolgen gilt zwar ebenfalls eine Ungleichung, aber eine andere und kompliziertere. Diese Unterscheidung ist prüfungsrelevant.
\end{merksatz}